% ============================================================
% Ejercicio 5
% ============================================================
\section{Ejercicio 5 (Opcional)}

\begin{tcolorbox}[enunciado]
Ejercicio opcional (obligatorio para los que trabajan en equipos de 4 integrantes)
\begin{itemize}[leftmargin=*]
    \item Describe e implementa un algoritmo que tome archivos del Ejercicio 2 y genere un archivo para el Ejercicio 4.
    \item Para resolver el problema general (cuando los pedidos tienen pesos diferentes), un alumno propone la siguiente idea: "Calcular el 'valor por gramo' (densidad) de cada pedido dividiendo su ganancia entre su peso. Ordenar los pedidos del más valioso por gramo al menos valioso. Meter los pedidos a la mochila en ese orden hasta que ya no quepa el siguiente. Si al terminar la ganancia es $\ge V$ responder SÍ".
\end{itemize}

¿Este enfoque siempre funciona de manera óptima? En caso de que no, diseña un ejemplar de prueba (un archivo) donde esta estrategia falle (es decir, el algoritmo diga que NO se alcanza la meta, pero en realidad SÍ existe una combinación válida), o justifica por qué este enfoque si podría siempre la respuesta correcta.
\end{tcolorbox}

\subsection{Solución}
% --- Escribe aquí tu solución ---